
\documentclass[11pt]{article}

\usepackage{amsmath}
\usepackage{amssymb}
\usepackage{geometry}
\usepackage{booktabs}
\usepackage{array}
\usepackage{hyperref}
\usepackage{enumitem}
\usepackage{caption}

\geometry{margin=1in}

\title{\textbf{Experiment Proposal: Optimal Batching Policies for ZK-Rollups Under Heterogeneous Transaction Loads}}
\author{}
\date{}

\begin{document}
\maketitle

\section{Problem Statement}

ZK-rollups amortize proof generation and L1 data availability costs by batching transactions. In practice, rollup workloads are highly heterogeneous: transactions differ substantially in proving complexity, execution cost, and data footprint. Current sequencer batching strategies (e.g., fixed-size or fixed-interval batching) are largely heuristic-driven and do not explicitly optimize for this heterogeneity.

\textbf{Core question:}
\begin{quote}
How should a centralized ZK-rollup sequencer batch heterogeneous transactions to minimize total proving and posting cost while satisfying latency constraints?
\end{quote}

This experiment provides a systematic, simulation-based evaluation of batching policies under realistic heterogeneous workloads, with the goal of producing actionable guidance for rollup designers and operators.

\section{Scope and Assumptions}

To ensure feasibility and clarity, the experiment is explicitly bounded:

\begin{itemize}[leftmargin=1.5em]
    \item Centralized sequencer with parallel provers (industry-standard model).
    \item Single rollup instance; no cross-rollup interactions.
    \item Trace-driven simulation (no live deployment).
    \item Evaluation of a small, controlled batching policy space.
\end{itemize}

\textbf{Out of scope:}
\begin{itemize}[leftmargin=1.5em]
    \item Decentralized sequencing
    \item MEV extraction strategies
    \item Cryptographic proof construction details
\end{itemize}

\section{Fee and Cost Model}

Rollup transaction fees reflect resource usage across both Layer 2 (off-chain execution and proving) and Layer 1 (data availability and verification).

\subsection{Total Transaction Fee}

\begin{equation}
Tx_{fee} = L2_{fee} + L1_{fee}
\end{equation}

\subsection{Layer 2 Fee Component}

\begin{equation}
L2_{fee} = (\rho_{l2} + \delta) \times g
\end{equation}

\begin{itemize}[leftmargin=1.5em]
    \item $\rho_{l2}$: L2 base fee (EIP-1559-style).
    \item $\delta$: Priority tip to the sequencer.
    \item $g$: L2 gas used for execution and amortized proving cost.
\end{itemize}

\subsection{Layer 1 Fee Component}

\begin{equation}
L1_{fee} = (\rho_{blob} \cdot Scalar_{blob} \cdot b_{tx}) + C_{tx} + S_{tx}
\end{equation}

\begin{itemize}[leftmargin=1.5em]
    \item $\rho_{blob}$: Blob base fee (or calldata equivalent).
    \item $b_{tx}$: Bytes contributed by the transaction.
    \item $C_{tx}$: L1 gas for batch commitment.
    \item $S_{tx}$: L1 gas for ZK proof verification.
\end{itemize}

The simulator focuses on \emph{relative cost comparisons} across batching policies rather than precise fee prediction.

\section{Transaction Cost Profiles}

Transactions are categorized into three abstract but realistic classes:

\begin{itemize}[leftmargin=1.5em]
    \item \textbf{Type A (Light)}: Simple transfers
    \begin{itemize}
        \item Low proving cost
        \item Low DA footprint
    \end{itemize}

    \item \textbf{Type B (Medium)}: ERC-20 transfers or swaps
    \begin{itemize}
        \item Moderate proving cost
        \item Moderate DA footprint
    \end{itemize}

    \item \textbf{Type C (Heavy)}: Complex smart contract calls
    \begin{itemize}
        \item High proving cost
        \item Relatively small DA footprint
    \end{itemize}
\end{itemize}

Cost parameters are derived from public rollup documentation, benchmarks, and academic literature, with sensitivity analysis applied to test robustness.

\section{Batching Policy Space}

\subsection{Baseline (Agnostic) Policies}

\begin{itemize}[leftmargin=1.5em]
    \item \textbf{Fixed-Size ($N$)}: Batch closes after $N$ transactions.
    \item \textbf{Fixed-Interval ($T$)}: Batch closes every $T$ seconds.
\end{itemize}

These policies ignore transaction heterogeneity and serve as control baselines.

\subsection{Soft-Cap (Reactive) Policies}

Used in production rollups (e.g., zkSync Era):
https://github.com/matter-labs/zksync-era/blob/main/core/node/state_keeper/src/seal_criteria/conditional_sealer.rs

\begin{quote}
Close the batch when the first resource limit is reached.
\end{quote}

\begin{itemize}[leftmargin=1.5em]
    \item Prover cycle budget
    \item L2 gas limit
    \item DA size limit
    \item Maximum latency timeout
\end{itemize}

\subsection{Transaction-Aware (Proactive) Policies}

\begin{itemize}[leftmargin=1.5em]
    \item \textbf{Isolation / Side Queues}: Heavy transactions scheduled only when prover slack exists.
    \item \textbf{Greedy Bin Packing}: Batch construction as a 2D packing problem
    \begin{itemize}
        \item Dimension 1: Proving cost
        \item Dimension 2: DA size
    \end{itemize}
    \item \textbf{Bounded Mixing}: Upper bounds on heavy transaction ratios per batch.
\end{itemize}

\section{Experimental Design}

\subsection{Workload Regimes}

\begin{itemize}[leftmargin=1.5em]
    \item Balanced steady-state workload
    \item Increasing skew toward heavy transactions
    \item Bursty arrivals under strict latency SLAs
\end{itemize}

\subsection{Experiment Phases}

\textbf{Phase 1: Baseline Benchmarking}
\begin{itemize}
    \item Steady workload
    \item Fixed-size and fixed-interval batching
\end{itemize}

\textbf{Phase 2: Heterogeneity Sensitivity}
\begin{itemize}
    \item Gradual increase in Type C fraction
    \item Constant arrival rate
\end{itemize}

\textbf{Phase 3: Latency-Constrained Evaluation}
\begin{itemize}
    \item SLA: max waiting time = 30s
    \item Compare soft-cap vs. transaction-aware policies
\end{itemize}

\section{Prover Performance Model}

\begin{itemize}[leftmargin=1.5em]
    \item Proving time modeled using abstract cycle counts.
    \item Batch proving latency determined by the slowest sub-proof plus recursion overhead.
    \item Parallel provers with finite capacity and queueing effects.
\end{itemize}

\section{Evaluation Metrics}

\begin{center}
\begin{tabular}{ll}
\toprule
\textbf{Metric} & \textbf{Purpose} \\
\midrule
Average / P99 latency & User experience and SLA compliance \\
Throughput (tx/s) & System efficiency \\
Cost per transaction & Economic efficiency \\
Prover utilization & Resource efficiency \\
Expired transactions & SLA violations \\
\bottomrule
\end{tabular}
\end{center}

\section{Hypotheses}

\begin{itemize}[leftmargin=1.5em]
    \item H1: Agnostic batching becomes unstable as transaction heterogeneity increases.
    \item H2: Soft-cap policies improve robustness but waste prover capacity.
    \item H3: Transaction-aware batching strictly dominates naive batching beyond a heterogeneity threshold.
\end{itemize}

\section{Expected Contributions}

\begin{itemize}[leftmargin=1.5em]
    \item An open-source, reusable batching simulator for ZK-rollups.
    \item Empirical identification of failure regimes for mixed batching.
    \item Clear cost–latency Pareto frontiers for batching policy selection.
    \item Practical guidance for sequencer and prover pipeline design.
\end{itemize}

\end{document}
```
